\section{Conclusiones}

Se sintetizaron AgTNPs obteniendo nanoplatos con morfología triangular y lados entre \qtyrange{20}{120}{\nano\meter}. Las imágenes de campo claro confirmaron la textura característica de estas estructuras, con las caras \hkl{111} depositadas preferentemente paralelas al sustrato de formvar, y su transparencia parcial al haz de electrones es un indicador indirecto del reducido espesor de los nanoplatos.

El análisis por transformada de Fourier de las imágenes de HRTEM permitió identificar reflexiones $\frac{1}{3}\{422\}$, prohibidas para la red fcc ideal de la plata pero características de regiones con estructura hcp. Esto confirma la presencia de fallas de apilamiento en los planos basales \hkl{111}, en concordancia con lo reportado en la literatura~\cite{rocha2007,tan2021}. El espaciado interplanar medido, \SI{0.26 \pm 0.1}{\nano\metre}, está en buen acuerdo con el valor de referencia de \SI{0.249}{\nano\metre}.

Las simulaciones multislice aportaron resultados cualitativos de relevancia directa. Una estructura fcc pura orientada con eje de zona \hkl[111] produce un contraste prácticamente nulo bajo las condiciones del Philips CM200-UT en el defoco de Scherzer, ya que el vector de onda correspondiente a los planos \hkl{220} cae en una región fuertemente atenuada por la CTF. Lo que indica de manera indirecta que toda AgTNP visible en una imagen de HRTEM debe presentar fallas de apilamiento. Por su parte, las estructuras puramente hcp presentan contraste apreciable que depende del espesor.

El estudio de estructuras tipo ``sánguche'' reveló que el contraste varía de manera no trivial incluso al mantener fija la cantidad de fallas en el centro, lo que indica que el contraste es una función compleja de la distribución global de fallas dentro de la partícula. Mostrando que el modelo de sánguche es insuficiente para un análisis detallado y es necesario explorar estructuras mas complejas. La metodología desarrollada en este trabajo sienta las bases para la construcción de una biblioteca de imágenes simuladas que, en el futuro, podría emplearse para entrenar una red neuronal capaz de identificar la distribución de fallas directamente a partir de imágenes experimentales.